\documentclass[10pt,a4paper]{amsart}
\usepackage{amsmath,amssymb,amsthm,graphicx}
\usepackage{verbatim}
\newenvironment{reference}{\comment}{\endcomment}
%\newenvironment{reference}{}{}
\newenvironment{slogan}{\comment}{\endcomment}
\newenvironment{history}{\comment}{\endcomment}

% For commutative diagrams we use Xy-pic
\usepackage[all]{xy}

% We use 2cell for 2-commutative diagrams.
\xyoption{2cell}
% Theorem environments.
%
\theoremstyle{plain}
\newtheorem{theorem}[subsection]{Theorem}
\newtheorem{proposition}[subsection]{Proposition}
\newtheorem{lemma}[subsection]{Lemma}

\theoremstyle{definition}
\newtheorem{definition}[subsection]{Definition}
\newtheorem{example}[subsection]{Example}
\newtheorem{exercise}[subsection]{Exercise}
\newtheorem{situation}[subsection]{Situation}

\theoremstyle{remark}
\newtheorem{remark}[subsection]{Remark}
\newtheorem{remarks}[subsection]{Remarks}

\numberwithin{equation}{subsection}

% Macros
%
\def\lim{\mathop{\mathrm{lim}}\nolimits}
\def\colim{\mathop{\mathrm{colim}}\nolimits}
\def\Spec{\mathop{\mathrm{Spec}}}
\def\Hom{\mathop{\mathrm{Hom}}\nolimits}
\def\Ext{\mathop{\mathrm{Ext}}\nolimits}
\def\SheafHom{\mathop{\mathcal{H}\!\mathit{om}}\nolimits}
\def\SheafExt{\mathop{\mathcal{E}\!\mathit{xt}}\nolimits}
\def\Sch{\mathit{Sch}}
\def\Mor{\mathop{\mathrm{Mor}}\nolimits}
\def\Ob{\mathop{\mathrm{Ob}}\nolimits}
\def\Sh{\mathop{\mathit{Sh}}\nolimits}
\def\NL{\mathop{N\!L}\nolimits}
\def\CH{\mathop{\mathrm{CH}}\nolimits}
\def\proetale{{pro\text{-}\acute{e}tale}}
\def\etale{{\acute{e}tale}}
\def\QCoh{\mathit{QCoh}}
\def\Ker{\mathop{\mathrm{Ker}}}
\def\Im{\mathop{\mathrm{Im}}}
\def\Coker{\mathop{\mathrm{Coker}}}
\def\Coim{\mathop{\mathrm{Coim}}}

% Boxtimes
%
\DeclareMathSymbol{\boxtimes}{\mathbin}{AMSa}{"02}

%
% Macros for moduli stacks/spaces
%
\def\QCohstack{\mathcal{QC}\!\mathit{oh}}
\def\Cohstack{\mathcal{C}\!\mathit{oh}}
\def\Spacesstack{\mathcal{S}\!\mathit{paces}}
\def\Quotfunctor{\mathrm{Quot}}
\def\Hilbfunctor{\mathrm{Hilb}}
\def\Curvesstack{\mathcal{C}\!\mathit{urves}}
\def\Polarizedstack{\mathcal{P}\!\mathit{olarized}}
\def\Complexesstack{\mathcal{C}\!\mathit{omplexes}}
% \Pic is the operator that assigns to X its picard group, usage \Pic(X)
% \Picardstack_{X/B} denotes the Picard stack of X over B
% \Picardfunctor_{X/B} denotes the Picard functor of X over B
\def\Pic{\mathop{\mathrm{Pic}}\nolimits}
\def\Picardstack{\mathcal{P}\!\mathit{ic}}
\def\Picardfunctor{\mathrm{Pic}}
\def\Deformationcategory{\mathcal{D}\!\mathit{ef}}
\usepackage{cleveref}
\crefname{originalproposition}{original proposition}{original propositions}
\Crefname{originalproposition}{Original proposition}{Original propositions}
\makeatletter
\newlabel{descent-proposition-equivalence-quasi-coherent-functorial}{{03LC}{source}{}{sourcepointer-1}{}}
\newlabel{descent-proposition-equivalence-quasi-coherent-functorial@cref}{{[originalproposition][0][]03LC}{[1][0][]source}{}{}{}}
\makeatother
\crefname{originalequation}{original equation}{original equations}
\Crefname{originalequation}{Original equation}{Original equations}
\makeatletter
\newlabel{equation-representable-pushforward}{{04CF}{source}{}{sourcepointer-2}{}}
\newlabel{equation-representable-pushforward@cref}{{[originalequation][0][]04CF}{[1][0][]source}{}{}{}}
\makeatother
\crefname{originallemma}{original lemma}{original lemmas}
\Crefname{originallemma}{Original lemma}{Original lemmas}
\makeatletter
\newlabel{lemma-base-change-quasi-compact}{{03HF}{source}{}{sourcepointer-3}{}}
\newlabel{lemma-base-change-quasi-compact@cref}{{[originallemma][0][]03HF}{[1][0][]source}{}{}{}}
\makeatother
\crefname{originallemma}{original lemma}{original lemmas}
\Crefname{originallemma}{Original lemma}{Original lemmas}
\makeatletter
\newlabel{lemma-base-change-separated}{{03KL}{source}{}{sourcepointer-4}{}}
\newlabel{lemma-base-change-separated@cref}{{[originallemma][0][]03KL}{[1][0][]source}{}{}{}}
\makeatother
\crefname{originallemma}{original lemma}{original lemmas}
\Crefname{originallemma}{Original lemma}{Original lemmas}
\makeatletter
\newlabel{lemma-compute-pushforward}{{03M8}{source}{}{sourcepointer-5}{}}
\newlabel{lemma-compute-pushforward@cref}{{[originallemma][0][]03M8}{[1][0][]source}{}{}{}}
\makeatother
\crefname{originallemma}{original lemma}{original lemmas}
\Crefname{originallemma}{Original lemma}{Original lemmas}
\makeatletter
\newlabel{spaces-properties-lemma-characterize-quasi-coherent}{{03M0}{source}{}{sourcepointer-6}{}}
\newlabel{spaces-properties-lemma-characterize-quasi-coherent@cref}{{[originallemma][0][]03M0}{[1][0][]source}{}{}{}}
\makeatother
\crefname{originallemma}{original lemma}{original lemmas}
\Crefname{originallemma}{Original lemma}{Original lemmas}
\makeatletter
\newlabel{spaces-properties-lemma-properties-quasi-coherent}{{03M1}{source}{}{sourcepointer-7}{}}
\newlabel{spaces-properties-lemma-properties-quasi-coherent@cref}{{[originallemma][0][]03M1}{[1][0][]source}{}{}{}}
\makeatother
\crefname{originallemma}{original lemma}{original lemmas}
\Crefname{originallemma}{Original lemma}{Original lemmas}
\makeatletter
\newlabel{spaces-properties-lemma-quasi-compact-affine-cover}{{03H6}{source}{}{sourcepointer-8}{}}
\newlabel{spaces-properties-lemma-quasi-compact-affine-cover@cref}{{[originallemma][0][]03H6}{[1][0][]source}{}{}{}}
\makeatother
\title{Pushforward of a quasi-coherent module along a quasi-compact quasi-separated morphism of algebraic spaces is quasi-coherent}
\author{The Stacks Project Authors}
\date{stacks-project:spaces-morphisms.tex:local-snapshot-20261012; GNU FDL 1.2 or later}
\begin{document}
\maketitle
\setcounter{subsection}{1}
\noindent\textit{Editorial note.} All delivered mathematical text is the author's own, copied from the pinned Stacks Project chapter file spaces-morphisms.tex; no statement or proof step was rewritten, reordered or removed.\par\medskip
\noindent\textit{Editorial note.} The counted proof uses equation (equation-representable-pushforward) of the same chapter as a labelled formula; the formula itself is consumed inside the proof span, so no extra context span was added.\par\medskip
\noindent\textit{Editorial note.} The eight references of the counted spans are recorded as non-load-bearing navigation pointers; the argument that consumes them stays inside the counted proof span.\par\medskip
\noindent\textit{Editorial note.} Preamble: the three native excerpt spans of sources/stacks-spaces-morphisms/preamble.tex (lines 8-13, 14-18 and 154-217). The chapter file itself carries no preamble text.\par\medskip
\noindent\textit{Allowed standard prerequisites (editorial).}\begin{itemize}
\item Quasi-coherence of a module on an algebraic space is tested on affine schemes etale over it. Properties of Spaces, Lemma spaces-properties-lemma-characterize-quasi-coherent, used at proof lines 2134-2138. Reduces the claim to a statement about the restriction of f\_*F to an affine etale V over Y.
\item The base change formula for pushforwards along an etale V -> Y. equation-representable-pushforward, used at proof lines 2144-2147. Identifies the restriction with the pushforward of the restriction along the base change f'.
\item The exact sequence computing the pushforward from an etale presentation. Lemma lemma-compute-pushforward, used at proof lines 2155-2159. Presents f\_*F as the kernel of a map between the pushforwards along the two projections of R.
\item Quasi-coherence of pushforwards of quasi-coherent modules under quasi-compact quasi-separated morphisms of schemes. Descent, Proposition descent-proposition-equivalence-quasi-coherent-functorial, used at proof lines 2164-2168. Makes a\_*(F|\_U) and b\_*(F|\_R) quasi-coherent.
\item A kernel of a map of quasi-coherent modules is quasi-coherent. Properties of Spaces, Lemma spaces-properties-lemma-properties-quasi-coherent, used at proof lines 2169-2172. Concludes for f\_*F.
\end{itemize}\medskip
\ifdefined\hypertarget\hypertarget{sourcepointer-1}{}\fi
\noindent\textit{Source reference.} Original proposition 03LC: focused-seeds/AGP/sources/stacks-snapshot-20261011/tags.txt, source lines 4708--4708. Navigation pointer into the Stacks Project; the standard fact it supplies is applied inside the counted proof span and the argument resting on it stays there.\par\medskip
\ifdefined\hypertarget\hypertarget{sourcepointer-2}{}\fi
\noindent\textit{Source reference.} Original equation 04CF: focused-seeds/AGP/sources/stacks-snapshot-20261011/tags.txt, source lines 5689--5689. Navigation pointer into the Stacks Project; the standard fact it supplies is applied inside the counted proof span and the argument resting on it stays there.\par\medskip
\ifdefined\hypertarget\hypertarget{sourcepointer-3}{}\fi
\noindent\textit{Source reference.} Original lemma 03HF: focused-seeds/AGP/sources/stacks-snapshot-20261011/tags.txt, source lines 4565--4565. Navigation pointer into the Stacks Project; the standard fact it supplies is applied inside the counted proof span and the argument resting on it stays there.\par\medskip
\ifdefined\hypertarget\hypertarget{sourcepointer-4}{}\fi
\noindent\textit{Source reference.} Original lemma 03KL: focused-seeds/AGP/sources/stacks-snapshot-20261011/tags.txt, source lines 4681--4681. Navigation pointer into the Stacks Project; the standard fact it supplies is applied inside the counted proof span and the argument resting on it stays there.\par\medskip
\ifdefined\hypertarget\hypertarget{sourcepointer-5}{}\fi
\noindent\textit{Source reference.} Original lemma 03M8: focused-seeds/AGP/sources/stacks-snapshot-20261011/tags.txt, source lines 4740--4740. Navigation pointer into the Stacks Project; the standard fact it supplies is applied inside the counted proof span and the argument resting on it stays there.\par\medskip
\ifdefined\hypertarget\hypertarget{sourcepointer-6}{}\fi
\noindent\textit{Source reference.} Original lemma 03M0: focused-seeds/AGP/sources/stacks-snapshot-20261011/tags.txt, source lines 4732--4732. Navigation pointer into the Stacks Project; the standard fact it supplies is applied inside the counted proof span and the argument resting on it stays there.\par\medskip
\ifdefined\hypertarget\hypertarget{sourcepointer-7}{}\fi
\noindent\textit{Source reference.} Original lemma 03M1: focused-seeds/AGP/sources/stacks-snapshot-20261011/tags.txt, source lines 4733--4733. Navigation pointer into the Stacks Project; the standard fact it supplies is applied inside the counted proof span and the argument resting on it stays there.\par\medskip
\ifdefined\hypertarget\hypertarget{sourcepointer-8}{}\fi
\noindent\textit{Source reference.} Original lemma 03H6: focused-seeds/AGP/sources/stacks-snapshot-20261011/tags.txt, source lines 4555--4555. Navigation pointer into the Stacks Project; the standard fact it supplies is applied inside the counted proof span and the argument resting on it stays there.\par\medskip
\begin{lemma}
\label{lemma-pushforward}
Let $S$ be a scheme.
Let $f : X \to Y$ be a morphism of algebraic spaces over $S$.
If $f$ is quasi-compact and quasi-separated, then $f_*$ transforms
quasi-coherent $\mathcal{O}_X$-modules into
quasi-coherent $\mathcal{O}_Y$-modules.
\end{lemma}

\begin{proof}
Let $\mathcal{F}$ be a quasi-coherent sheaf on $X$. We have to show that
$f_*\mathcal{F}$ is a quasi-coherent sheaf on $Y$. For this it suffices
to show that for any affine scheme $V$ and \'etale morphism $V \to Y$
the restriction of $f_*\mathcal{F}$ to $V$ is quasi-coherent, see
Properties of Spaces,
Lemma \ref{spaces-properties-lemma-characterize-quasi-coherent}.
Let $f' : V \times_Y X \to V$
be the base change of $f$ by $V \to Y$. Note that $f'$ is also
quasi-compact and quasi-separated, see
Lemmas \ref{lemma-base-change-quasi-compact} and
\ref{lemma-base-change-separated}.
By (\ref{equation-representable-pushforward})
we know that the restriction of $f_*\mathcal{F}$ to $V$ is $f'_*$ of the
restriction of $\mathcal{F}$ to $V \times_Y X$. Hence
we may replace $f$ by $f'$, and assume that $Y$ is an affine scheme.

\medskip\noindent
Assume $Y$ is an affine scheme. Since $f$ is quasi-compact we see that $X$
is quasi-compact. Thus we may choose an affine scheme $U$ and a surjective
\'etale morphism $U \to X$, see
Properties of Spaces,
Lemma \ref{spaces-properties-lemma-quasi-compact-affine-cover}.
By Lemma \ref{lemma-compute-pushforward} we get an exact sequence
$$
0 \to f_*\mathcal{F} \to a_*(\mathcal{F}|_U) \to b_*(\mathcal{F}|_R).
$$
where $R = U \times_X U$.
As $X \to Y$ is quasi-separated we see that $R \to U \times_Y U$
is a quasi-compact monomorphism. This implies that $R$ is a quasi-compact
separated scheme (as $U$ and $Y$ are affine at this point).
Hence $a : U \to Y$ and $b : R \to Y$ are quasi-compact and
quasi-separated morphisms of schemes. Thus by
Descent,
Proposition \ref{descent-proposition-equivalence-quasi-coherent-functorial}
the sheaves $a_*(\mathcal{F}|_U)$ and $b_*(\mathcal{F}|_R)$
are quasi-coherent (see also the discussion preceding this lemma).
This implies that $f_*\mathcal{F}$ is a kernel of
quasi-coherent modules, and hence itself quasi-coherent, see
Properties of Spaces,
Lemma \ref{spaces-properties-lemma-properties-quasi-coherent}.
\end{proof}

\end{document}
